\documentclass[../../../../main.tex]{subfiles}
\begin{document}

\begin{flushright}
\footnotesize\textit{【自動文字起こし・要確認】}
\end{flushright}

\noindent
[\textbf{解}] 1チームの勝ち数は、$n-1, n-2, \dots, 1, 0$ の $n$ 通りしかないので、題意から、第 $k$ 位のチーム $A_k$ は $n-k$ 勝していることになる ($k=1, 2, \dots, n$)。

\noindent
以下、題意の成立を $k$ についての帰納法で示す。

\begin{enumerate}
    \item[$1^\circ$] $k=1$ の時\\
    1位のチームは $n-1$ 勝、つまり他の全てのチームに勝っており、成立。

    \item[$2^\circ$] $k=i$ での成立を仮定 ($i=1, 2, \dots, n-1$)\\
    $k=i+1$ の時、仮定から、$A_{i+1}$ は $A_1 \sim A_i$ のチームに負けており、少なくとも $i$ 回負けている。従って、$n-(i+1)$ 勝するには残りの全ての試合に勝つことになるから、$k=i+1$ でも成立。
\end{enumerate}

\noindent
以上から示された。

\newpage

\end{document}
